
\subsection*{A6: Conversational Chess}

 
Using \libName, we developed Conversational Chess, a natural language interface game shown in the last sub-figure of Figure~\ref{fig:app_summary}.
 It features built-in agents for players, which can be human or LLM, and a third-party board agent to provide information and validate moves based on standard rules.

With \libName, we enabled two essential features: (1) Natural, flexible, and engaging game dynamics, enabled by the customizable agent design in \libName. Conversational Chess supports a range of game-play patterns, including AI-AI, AI-human, and human-human, with seamless switching between these modes during a single game. An illustrative example of these entertaining game dynamics can be found in Figure~\ref{fig:conversational-chess-example}, Appendix~\ref{appendix:app:chess}. (2) Grounding, which is a crucial aspect to maintain game integrity. During gameplay, the board agent checks each proposed move for legality; if a move is invalid, the agent responds with an error, prompting the player agent to re-propose a legal move before continuing. This process ensures that only valid moves are played and helps maintain a consistent gaming experience.
As an ablation study, we removed the board agent and instead only relied on a relevant prompt \textit{``you should make sure both you and the opponent are making
legal moves"} to ground their move. 
The results highlighted that without the board agent, illegitimate moves caused game disruptions. 
The modular design offered flexibility, allowing swift adjustments to the board agent in response to evolving game rules or varying chess rule variants. A comprehensive demonstration of this ablation study is in Appendix~\ref{appendix:app:chess}.


